\documentclass[12pt]{article}
\usepackage{graphicx}
\usepackage[T1]{fontenc}
\usepackage{tocloft}
\usepackage{todonotes}
\usepackage{subcaption}
\usepackage{caption}
\usepackage{hyperref}
\usepackage{booktabs}
\usepackage{listings}
\usepackage{wrapfig}
\usepackage{pdfpages}
\usepackage{pdflscape}
\usepackage{scrlayer}
\usepackage{scrhack}
\begin{document}
\begin{titlepage}
    \centering
    {\scshape\large AA 2020/2021 \par}
    \vfill
    \includegraphics[width=100pt]{assets/logo-polimi-new}\par\vspace{1cm}
    {\scshape\LARGE Politecnico di Milano \par}
    \vspace{1.5cm}
    {\huge\bfseries RASD\@: Requirement Analysis
        and Specification Document \par}
    \vspace{2cm}
    {\Large {Alice Piemonti\quad Luca Pirovano\quad Nicolò Sonnino}\par}
    \vfill
    {\large Professor\par
        Matteo \textsc{Rossi}}
    \vfill
    {\large \today \par}
\end{titlepage}
\tableofcontents
\newpage
\section{Introduction}
CLup (Customers Line-up) is an easy-to-use application which intent is to limit as much as possible the number of accesses to grocery shopping.\\

Due to the recent worldwide spread of SARS-CoV-2 (COVID-19), many countries are imposing lockdowns and strict regulations about social distancing such as: the closure of restaurants in the evening, limitations on public transports, curfews, etc.

In particular, citizens are experiences difficulties in accessing supermarkets because of their limited capacity.
\subsection{Purpose}
The aim of the product is to avoid gatherings outside grocery stores, improving the safety of the customers.

This is achieved through monitoring accesses to the buildings, managing time slots for visits and optimizing people flows inside the stores.\\

The application will be operable freely, widely available and very intuitive because \todo{completare}

The userbase is expected to be both people with an Internet access and ones without it, from young people to elderly.

\subsection{Scope}
The product shall be called CLup and will let users to plan their shopping session in two different ways:
\begin{itemize}
    \item {\textbf{ASAP}: the user will claim the first available ticket and receive an estimated queue time.}
    \item {\textbf{Reservation}: the user will choose a time slot from a list of available ones, in order to book his visit to the structure.}
\end{itemize}

Every customer can choose one of these modes \textbf{remotely} via an official app or through a web browser, or \textbf{in presence} by asking to a staff member, who will act as an intermediate between the customer and the system.

When a customer makes a reservation, the system allows him to choose the duration of his visit and insert a list of possible purchases, in order to optimize his stay.\\

In addition, the user can change time slot/store based on system's suggestions and enable periodically notifications of available slots in a day/time range.


\newpage
\section{Overall Description}
\newpage
\section{Specific Requirements}
\newpage
\section{Formal Analysis}
\newpage
\section{Effort spent}
\end{document}